% 図9-3 変数の型とサイズ(2)  x23_modeling4.tex から \input する
\begin{tikzpicture}[
  head/.style={font=\large},
  decl/.style={font=\ttfamily},
  dim/.style={font=\small}]
  \colorlet{stanblue}{blue!60!cyan}
  \def\xa{1.6}\def\xb{5.8}\def\xc{11.0}
  % 長さKの縦棒: \kbar{色}{左下x}{左下y}
  \newcommand{\kbar}[3]{\filldraw[fill=#1, draw=white, line width=0.6pt] (#2,#3) rectangle ++(0.36,2.2);}
  % 見出し
  \node[head] at (\xa,0) {型};
  \node[head] at (\xb,0) {宣言例};
  \node[head] at (\xc,0) {サイズ};
  \draw[line width=1pt] (0,-0.45) -- (13.6,-0.45);
  % ベクトル1本
  \node at (\xa,-1.0) {ベクトル};
  \node[decl] at (\xb,-1.0) {vector[K] X};
  \kbar{stanblue}{\xc-0.18}{-3.0}
  \node[dim, left] at (\xc-0.2,-2.85) {K};
  \draw[line width=0.4pt] (0,-3.3) -- (13.6,-3.3);
  % ベクトルがN本
  \node at (\xa,-3.85) {ベクトルの配列};
  \node[decl] at (\xb,-3.85) {array[N] vector[K] X};
  \foreach \x in {0,0.5,1.0,2.3} \kbar{stanblue}{\xc-1.3+\x}{-6.1}
  \node at (\xc+0.42,-5.0) {$\cdots$};
  \node[dim, left] at (\xc-1.32,-5.95) {K};
  \node[dim, above] at (\xc+1.18,-3.9) {N};
  \draw[line width=0.4pt] (0,-6.45) -- (13.6,-6.45);
  % ベクトルがN×M本(奥行き方向にM)。奥の層から順に描いて手前で隠す
  \node at (\xa,-7.0) {ベクトルの2次元配列};
  \node[decl] at (\xb,-7.0) {array[N,M] vector[K] X};
  \foreach \x in {0,0.5,1.0,2.3} \kbar{stanblue!25}{10.3+\x}{-9.65}
  \foreach \x in {0,0.5,1.0,2.3} \kbar{stanblue!45}{9.9+\x}{-10.1}
  \foreach \x in {0,0.5,1.0,2.3} \kbar{stanblue!70}{9.5+\x}{-10.55}
  \foreach \x in {0,0.5,1.0,2.3} \kbar{stanblue!100}{9.1+\x}{-11.0}
  \node at (10.82,-9.9) {$\cdots$};
  \node[dim, left] at (9.08,-10.85) {K};
  \node[dim, above] at (12.78,-7.45) {N};
  \draw[-{Stealth[length=3mm]}, line width=0.8pt] (12.0,-11.15) -- (13.2,-9.8)
    node[dim, right] {M};
  \draw[line width=1pt] (0,-11.4) -- (13.6,-11.4);
\end{tikzpicture}
